\documentclass{situationist-preprint}

\title{Remorse Without Wrongdoing}
\msid{ETH-27-b}
\msfamily{ETH-27}
\msdiscipline{Normative and Applied Ethics}
\mstype{Companion manuscript}
\mskeywords{remorse; infringement; dirty hands; agent-regret; fitting emotions; compensation}

\begin{document}
\maketitle

\begin{abstract}
ETH-27-a argues that the politician who dirties her hands for a greater good acts rightly and that remorse is nonetheless fitting, because remorse presents its object as the infringement of a serious claim, not as a wrong. We show that the structure is not peculiar to politics. Whenever an agent permissibly infringes a serious claim of another, remorse fits her act and guilt does not. We illustrate with a breach of property to save one's life, a broken promise to meet an emergency, and a forced choice between the claims of two persons, and we distinguish the resulting remorse from Williams's agent-regret, which fits outcomes caused through one's agency without the infringement of any claim.
\end{abstract}

\section{The general structure}

ETH-27-a defines remorse as a painful, backward-looking emotion that presents one's own act as having infringed, through one's own agency, a serious claim of another (its Definition~2.3), and distinguishes infringing a claim from violating it, that is, infringing it impermissibly (Definition~2.1, after \citealt[ch.~3]{Thomson1990}). By its Principle~P1, an emotion is fitting if and only if what it presents is the case. The argument of its Theorem~3.2 uses no feature of politics.

\begin{proposition}\label{p:general}
Whenever an agent permissibly infringes a serious claim of another, remorse for the act is fitting, and guilt, understood as presenting the act as wrong, is not.
\end{proposition}

\begin{proof}
What remorse presents, the infringement of a serious claim through one's agency, is the case; so remorse is fitting by Principle~P1 of ETH-27-a. What guilt presents, that the act was wrong, is not the case, since the infringement was permissible; so guilt is not fitting.
\end{proof}

Proposition~\ref{p:general} is the form the conclusion of ETH-27 takes once the content of remorse is fixed, and it fits the view of MND-17 that emotions represent evaluative properties: remorse represents an act as an infringement, which is an evaluative property distinct from wrongness.

\section{Cases}

\paragraph{The cabin.} A hiker caught in a blizzard breaks into an unoccupied cabin, burns the owner's furniture for warmth and eats his food \citep{Feinberg1978}. She acts permissibly, and she infringes the owner's claims. She owes him acknowledgment and compensation, as Proposition~3.4 of ETH-27-a predicts, and remorse, directed at what she did to his property, fits her situation in a way that guilt does not.

\paragraph{The broken promise.} An agent who has promised to attend a friend's ceremony instead takes a stranger injured in the street to hospital. Breaking the promise is permissible, but the friend's claim, created by the promise, is infringed. An apology to the friend is due, and its sincerity requires the attitude that remorse is: not that one acted wrongly, but that one failed someone who had a claim on one.

\paragraph{The forced choice.} Where an agent can honor the claim of only one of two persons, ETH-15 holds that one option is uniquely permissible, and the agent who takes it acts rightly. The other person's claim is infringed, and remorse toward that person is fitting. This is the structure of the case examined in ETH-15-b, and it shows how the denial of genuine moral dilemmas is consistent with the moral residue that dilemma theorists rightly insist on.

\section{Remorse and agent-regret}

\citet{Williams1976} described the lorry driver who, through no fault of his own, runs over a child, and argued that he fittingly feels a regret that a spectator does not share. If the driver was faultless, he infringed no claim, since no one has a claim against others that they not be the faultless causes of harm. Agent-regret then fits him, since the harm came about through his agency, but remorse does not. The two emotions are distinguished by what they present: agent-regret, an outcome brought about through one's agency; remorse, the infringement of a claim. Dirty hands, and the cases above, call for remorse because a claim is infringed, which is why amends in these cases are owed rather than merely offered.

\begin{objection}
The distinction between remorse and guilt is too fine to be felt; an emotion that presents one's act as an infringement of another's serious claim will be experienced as guilt.
\end{objection}

\begin{reply}
Fittingness concerns content, not phenomenology \citep{DArmsJacobson2000}. Whether or not the two feel alike, they differ in what they present and in what they dispose one to. Remorse disposes to acknowledgment and amends toward the person whose claim was infringed; guilt, presenting the act as wrong, disposes in addition to self-reproach and to accepting blame and sanction as deserved, which in a case of permissible infringement would be a mistake.
\end{reply}

\section{Conclusion}

Remorse without wrongdoing is the ordinary response to the permissible infringement of a claim, in private life as in politics. The dirty-hands politician is a dramatic case of a common structure.

\begin{thebibliography}{9}\small
\bibitem[D'Arms and Jacobson(2000)]{DArmsJacobson2000} D'Arms, J. and Jacobson, D. (2000). The moralistic fallacy: On the `appropriateness' of emotions. \emph{Philosophy and Phenomenological Research}, 61(1), 65--90.
\bibitem[Feinberg(1978)]{Feinberg1978} Feinberg, J. (1978). Voluntary euthanasia and the inalienable right to life. \emph{Philosophy \& Public Affairs}, 7(2), 93--123.
\bibitem[Thomson(1990)]{Thomson1990} Thomson, J.~J. (1990). \emph{The Realm of Rights}. Harvard University Press.
\bibitem[Williams(1976)]{Williams1976} Williams, B. (1976). Moral luck. \emph{Proceedings of the Aristotelian Society, Supplementary Volumes}, 50, 115--135.
\end{thebibliography}

\end{document}
